\chapter{La théorie des groupes en action}\label{ch:b3:group-theory-applied}

Le dioxyde de carbone ne représente que quelques centaines de molécules
par million dans l'air, le diazote et le dioxygène presque tout le
reste~; pourtant c'est le dioxyde de carbone, et non le diazote ou le
dioxygène, qui absorbe le rayonnement infrarouge que le sol chaud envoie
vers l'espace. Une molécule n'absorbe l'infrarouge que par une vibration
qui fait varier son moment dipolaire, et c'est la seule symétrie qui
décide si une vibration le fait~: l'unique vibration du diazote ne le
peut pas, deux des quatre vibrations du dioxyde de carbone le peuvent. Ce
chapitre fait des \omterm{def:b3:point-groups:irrep}{tables de caractères} du \cref{ch:b3:point-groups} des
outils~: il réduit les représentations, construit des orbitales adaptées
à la symétrie, dénombre et nomme les vibrations, et établit les règles de
sélection des spectroscopies infrarouge et Raman.

\begin{recall}
Le \Cref{ch:b3:point-groups} a défini les \omterm{def:b3:point-groups:operation}{opérations de symétrie}, les
\omterm{def:b3:point-groups:point-group}{groupes ponctuels}, les classes, les représentations, les \omterm{def:b3:point-groups:representation}{caractères} et
les \omterm{def:b3:point-groups:irrep}{tables de caractères}. Le volume de deuxième année a construit les
orbitales moléculaires de \ce{H2O}, \ce{NH3} et \ce{CH4} à partir
d'orbitales de fragment et de combinaisons des orbitales $1s$ des
hydrogènes adaptées à la symétrie, en les désignant par des symboles
établis ici. Le volume de première année a lu des spectres infrarouges
en nombres d'onde et défini la polarisabilité d'une molécule.
\end{recall}

\section{Réduire une représentation}

Tout ensemble de fonctions ou de vecteurs attaché à une molécule porte une
représentation de son \omterm{def:b3:point-groups:point-group}{groupe ponctuel}. La plupart sont réductibles.

\begin{definition}[Représentation réductible, représentation totalement symétrique]\label{def:b3:group-theory-applied:reducible}
Une \emph{représentation réductible}\index{representation reductible@représentation réductible} est une
représentation dont toutes les matrices peuvent être mises, par un même
changement de base, sous la même forme diagonale par blocs~; c'est alors
une somme de \omterm{def:b3:point-groups:irrep}{représentations irréductibles},
notée $\Gamma = \sum_in_i\Gamma_i$. La \omterm{def:b3:point-groups:irrep}{représentation irréductible} dont
tous les \omterm{def:b3:point-groups:representation}{caractères} valent 1 est la \emph{représentation totalement
symétrique}\index{representation totalement symetrique@représentation totalement symétrique} ($A_1$, $A_{1g}$,
$A_1'$\dots).
\end{definition}

\begin{theorem}[Grand théorème d'orthogonalité]\label{thm:b3:group-theory-applied:great-orthogonality}
Pour deux \omterm{def:b3:point-groups:irrep}{représentations irréductibles} $\Gamma_i$, $\Gamma_j$ de
dimensions $d_i$,
$d_j$ d'un groupe d'ordre $h$, écrites avec des matrices unitaires,
\[
\sum_R D_i(R)_{mn}^*\,D_j(R)_{m'n'} = \frac{h}{d_i}\,\delta_{ij}\,\delta_{mm'}\,\delta_{nn'}.
\]
\end{theorem}
\admitted

La démonstration, par les lemmes de Schur, relève de cours plus avancés~;
ce sont ses conséquences pour les \omterm{def:b3:point-groups:representation}{caractères} que la chimie utilise, et
chaque \omterm{def:b3:point-groups:irrep}{table de caractères} de ce livre est vérifiée à leur aune.

\begin{corollary}[Orthogonalité des caractères]\label{cor:b3:group-theory-applied:characters}
En notant $g_c$ le nombre d'opérations de la classe $c$,
\[
\sum_cg_c\,\chi_i(c)^*\chi_j(c) = h\,\delta_{ij},
\]
et le nombre de \omterm{def:b3:point-groups:irrep}{représentations irréductibles} est égal au nombre de
classes,
avec $\sum_id_i^2 = h$.
\end{corollary}

\begin{proof}
Posons $m = n$ et $m' = n'$ dans le théorème et sommons sur $m$ et $m'$~:
le membre de gauche devient $\sum_R\chi_i(R)^*\chi_j(R)$, celui de droite $\frac{h}{d_i}\delta_{ij}
\sum_{m}\delta_{mm} = h\delta_{ij}$~; regroupons les opérations par
classe. Les lignes de la table sont donc des vecteurs orthogonaux d'un
espace dont la dimension est le nombre de classes, de sorte qu'il y a au
plus autant de \omterm{def:b3:point-groups:irrep}{représentations irréductibles} que de classes~; l'égalité
et $\sum d_i^2 = h$ s'obtiennent en appliquant le théorème à la
représentation régulière (exercice 10).
\end{proof}

\begin{theorem}[La formule de réduction]\label{thm:b3:group-theory-applied:reduction}
Une représentation $\Gamma$ de \omterm{def:b3:point-groups:representation}{caractères} $\chi(c)$ contient la
\omterm{def:b3:point-groups:irrep}{représentation irréductible} $\Gamma_i$ exactement
\[
n_i = \frac1h\sum_cg_c\,\chi(c)\,\chi_i(c)^*
\]
fois. C'est la \emph{formule de réduction}\index{formule de reduction@formule de réduction}.
\end{theorem}

\begin{proof}
Comme $\Gamma = \sum_jn_j\Gamma_j$, ses \omterm{def:b3:point-groups:representation}{caractères} sont $\chi(c) = \sum_jn_j\chi_j(c)$.
Multiplions par $g_c\chi_i(c)^*$, sommons sur les classes et utilisons le
corollaire~:
$\sum_cg_c\chi(c)\chi_i(c)^* = \sum_jn_jh\delta_{ij} = hn_i$.
\end{proof}

\begin{example}[Les orbitales des hydrogènes de l'ammoniac]\label{ex:b3:group-theory-applied:nh3}
Prenons pour base les trois orbitales $1s$ des hydrogènes de \ce{NH3}.
$E$ les laisse toutes trois en place ($\chi = 3$), $C_3$ les déplace
toutes trois ($\chi = 0$), un
$\sigma_v$ en laisse une en place et échange les deux autres
($\chi = 1$). Avec la table de $C_{3v}$~:
$n_{A_1} = \frac16(3 + 0 + 3\times1) = 1$, $n_{A_2} = \frac16(3 + 0 - 3) = 0$,
$n_E = \frac16(6 + 0 + 0) = 1$. Donc $\Gamma_{\mathrm H} = A_1 + E$~: les
trois hydrogènes donnent une combinaison totalement symétrique et une
paire dégénérée.
\end{example}

\begin{method}[Réduire une représentation]\label{met:b3:group-theory-applied:reduce}
\begin{enumerate}
  \item Pour chaque classe, trouver le \omterm{def:b3:point-groups:representation}{caractère}~: pour un ensemble de
        fonctions centrées sur les atomes, le nombre de fonctions laissées
        en place, chacune comptée avec le signe qu'elle acquiert.
  \item Appliquer la \omterm{thm:b3:group-theory-applied:reduction}{formule de réduction} avec la ligne de chaque
        \omterm{def:b3:point-groups:irrep}{représentation irréductible}, en pondérant chaque classe par son
        effectif.
  \item Vérifier~: les résultats sont des entiers positifs ou nuls, et
        $\sum_in_id_i$ est égal au \omterm{def:b3:point-groups:representation}{caractère} sous $E$ (la taille de la
        base).
\end{enumerate}
\end{method}

\section{Orbitales adaptées à la symétrie}

\begin{definition}[Opérateur de projection]\label{def:b3:group-theory-applied:projection}
L'\emph{opérateur de projection}\index{operateur de projection@opérateur de projection} sur la
\omterm{def:b3:point-groups:irrep}{représentation irréductible} $\Gamma_i$ est
\[
\hat P_i = \frac{d_i}{h}\sum_R\chi_i(R)^*\,\hat R,
\]
où $\hat R$ applique l'opération $R$ à une fonction.
\end{definition}

\begin{proposition}[Ce que fait l'opérateur de projection]\label{prop:b3:group-theory-applied:projection}
Appliqué à une fonction quelconque $f$, $\hat P_if$ est soit nul, soit une
fonction qui se transforme comme $\Gamma_i$ (une combinaison adaptée à
la symétrie)~; pour une représentation $\Gamma_i$ de dimension un, elle
est inchangée, au signe $\chi_i(S)$ près, par chaque
opération $S$.
\end{proposition}

\begin{proof}[Démonstration partielle]
Pour $d_i = 1$~: $\hat S\hat P_if = \frac1h\sum_R\chi_i(R)\hat S\hat Rf$.
Posons $R' = SR$,
qui parcourt tout le groupe quand $R$ le parcourt~; $\chi_i(R) = \chi_i(S^{-1}R') =
\chi_i(S)\chi_i(R')$ pour une représentation de dimension un (les
\omterm{def:b3:point-groups:representation}{caractères} sont les matrices elles-mêmes). Donc $\hat S\hat P_if = \chi_i(S)\hat P_if$.
Le cas général utilise le grand théorème d'orthogonalité et est admis.
\end{proof}

\begin{method}[Construire des combinaisons adaptées à la symétrie]\label{met:b3:group-theory-applied:salc}
\begin{enumerate}
  \item Choisir une fonction de l'ensemble, $h_1$, et noter dans un
        tableau la fonction sur laquelle chaque opération l'envoie.
  \item Pour chaque \omterm{def:b3:point-groups:irrep}{représentation irréductible}, multiplier chaque image
        par le \omterm{def:b3:point-groups:representation}{caractère} de l'opération, additionner et normer (en
        négligeant le recouvrement entre fonctions atomiques).
  \item Pour une représentation dégénérée, projeter une seconde fonction
        pour obtenir un second membre, puis orthogonaliser.
\end{enumerate}
\end{method}

\begin{example}[Les combinaisons de l'ammoniac et de l'eau]\label{ex:b3:group-theory-applied:salcs}
Dans \ce{NH3}, $\hat P_{A_1}h_1 \propto h_1 + h_2 + h_3$ (tous les
\omterm{def:b3:point-groups:representation}{caractères} valent 1). Pour $E$
(\omterm{def:b3:point-groups:representation}{caractères} $2, -1, 0$), $\hat P_Eh_1 \propto 2h_1 - h_2 - h_3$~; projeter
$h_2$ et orthogonaliser donne le partenaire $h_2 - h_3$. Après
normalisation~: $a_1 =
\frac{1}{\sqrt3}(h_1 + h_2 + h_3)$, $e = \frac{1}{\sqrt6}(2h_1 - h_2 - h_3)$,
$\frac{1}{\sqrt2}(h_2 - h_3)$. Dans \ce{H2O}, placée dans le plan $xz$,
$h_1 + h_2$ est
$a_1$ et $h_1 - h_2$ est $b_1$ (elle change de signe par $C_2$ et par
$\sigma_v'(yz)$, qui échangent les hydrogènes). Avec la molécule dans le
plan $yz$, comme le préfèrent certains ouvrages, les symboles $b_1$ et
$b_2$ sont échangés partout.
\end{example}

\begin{omfigure}
\begin{tikzpicture}[font=\small]
  \foreach \x/\lab/\ca/\cb/\cc in {0/{$a_1$}/1/1/1, 4/{$e$}/2/-1/-1, 8/{$e$}/0/1/-1}{
    \begin{scope}[xshift=\x cm]
      \foreach \a/\c in {90/\ca,210/\cb,330/\cc}{
        \pgfmathsetmacro{\r}{0.18+0.13*abs(\c)}
        \ifdim \c pt > 0pt \fill[omDef!55, draw=omDef] (\a:1.0) circle (\r);
        \else \ifdim \c pt < 0pt \fill[omThm!25, draw=omThm] (\a:1.0) circle (\r);
        \else \draw[gray, densely dotted] (\a:1.0) circle (0.12); \fi\fi
      }
      \node[aN, minimum size=4.6mm] at (0,0) {};
      \node[anchor=north] at (0,-1.45) {\lab};
    \end{scope}}
  \node[anchor=north, font=\footnotesize] at (0,-1.85) {$h_1 + h_2 + h_3$};
  \node[anchor=north, font=\footnotesize] at (4,-1.85) {$2h_1 - h_2 - h_3$};
  \node[anchor=north, font=\footnotesize] at (8,-1.85) {$h_2 - h_3$};
\end{tikzpicture}

{\small Les combinaisons adaptées à la symétrie des trois orbitales $1s$
des hydrogènes de l'ammoniac, vues selon l'axe $C_3$ ($h_1$ en haut).
Bleu~: coefficient positif~; rouge~: négatif, la taille de chaque disque
croissant avec le coefficient~; pointillés~: nul.}
\end{omfigure}

\begin{example}[Six ligands autour d'un métal]\label{ex:b3:group-theory-applied:ml6}
Pour six orbitales $\sigma$ de ligands pointant vers un métal depuis les
sommets d'un octaèdre, les \omterm{def:b3:point-groups:representation}{caractères} sur les classes de $O_h$ ($E$,
$8C_3$, $6C_2$, $6C_4$,
$3C_2$, $i$, $6S_4$, $8S_6$, $3\sigma_h$, $6\sigma_d$) sont $6, 0, 0, 2, 2, 0, 0, 0, 4,
2$, et la réduction donne $\Gamma_\sigma = A_{1g} + E_g + T_{1u}$. Les
orbitales $s$ du métal
($a_{1g}$), $d_{z^2}$ et $d_{x^2-y^2}$ ($e_g$) et $p$ ($t_{1u}$) trouvent
des partenaires parmi les ligands~; ses orbitales $d_{xy}$, $d_{xz}$,
$d_{yz}$ ($t_{2g}$) n'en trouvent aucun et restent non liantes
en $\sigma$ --- c'est l'origine de l'éclatement $t_{2g}$/$e_g$ du volume
de deuxième année, développé aux \cref{ch:b3:organometallic-bonding,ch:b3:complex-spectra-magnetism}.
\end{example}

\begin{omfigure}
\begin{tikzpicture}[font=\small, yscale=0.95]
  % O orbitals (left), MOs (centre), H combinations (right); schematic energies
  \foreach \y/\l in {0/{$2s\ (a_1)$}}{\draw[omstate] (0,\y) -- (1.2,\y); \node[anchor=east] at (-0.05,\y) {\l};}
  \draw[omstate] (0,3.0) -- (0.36,3.0) (0.42,3.0) -- (0.78,3.0) (0.84,3.0) -- (1.2,3.0);
  \node[anchor=east] at (-0.05,3.0) {$2p\ (b_1, a_1, b_2)$};
  \draw[omstate] (6.8,3.9) -- (8.0,3.9); \node[anchor=west] at (8.05,3.9) {$h_1 - h_2\ (b_1)$};
  \draw[omstate] (6.8,3.6) -- (8.0,3.6); \node[anchor=west] at (8.05,3.5) {$h_1 + h_2\ (a_1)$};
  \foreach \y/\l in {-0.6/{$2a_1$}, 1.6/{$1b_1$}, 2.4/{$3a_1$}, 3.0/{$1b_2$}, 5.6/{$4a_1$}, 6.2/{$2b_1$}}{
    \draw[omstate] (3.4,\y) -- (4.6,\y); \node[anchor=west] at (4.65,\y) {\l};}
  \foreach \y in {-0.6,1.6,2.4,3.0}{\node[font=\scriptsize] at (4.0,\y+0.16) {$\uparrow\downarrow$};}
  \draw[densely dotted, gray] (1.2,0) -- (3.4,-0.6) (1.2,3.0) -- (3.4,1.6) (1.2,3.0) -- (3.4,2.4) (1.2,3.0) -- (3.4,3.0)
     (6.8,3.6) -- (4.6,-0.6) (6.8,3.9) -- (4.6,1.6) (6.8,3.6) -- (4.6,2.4) (1.2,3.0) -- (3.4,6.2) (6.8,3.9) -- (4.6,6.2) (6.8,3.6) -- (4.6,5.6);
  \node[anchor=north] at (0.6,-1.0) {O};
  \node[anchor=north] at (4.0,-1.0) {\ce{H2O}};
  \node[anchor=north] at (7.4,-1.0) {2 H};
  \draw[->] (-2.3,-0.9) -- (-2.3,6.4) node[above] {$E$};
\end{tikzpicture}

{\small Orbitales moléculaires de valence de l'eau (énergies
schématiques, molécule dans le plan $xz$). Seules des orbitales de même
symétrie se mélangent~: la combinaison $b_1$ des hydrogènes avec $2p_x$,
la combinaison $a_1$ avec $2s$ et $2p_z$~; $2p_y$ ($b_2$) n'a pas de
partenaire et reste un doublet non liant, l'orbitale occupée la plus
haute.}
\end{omfigure}

\section{Modes de vibration}

\begin{definition}[Mode normal]\label{def:b3:group-theory-applied:normal-mode}
Un \emph{mode normal}\index{mode normal} d'une molécule est une vibration
collective dans laquelle tous les atomes oscillent à la même fréquence et
en phase, selon un vecteur propre de la hessienne pondérée par les masses
(\cref{prop:b3:computational-chemistry:hessian}). Chaque mode normal se
transforme selon une \omterm{def:b3:point-groups:irrep}{représentation irréductible} du \omterm{def:b3:point-groups:point-group}{groupe ponctuel}.
\end{definition}

\begin{proposition}[La représentation de tous les mouvements]\label{prop:b3:group-theory-applied:gamma-3n}
Attachons trois vecteurs déplacement $x$, $y$, $z$ à chaque atome. Pour
une opération
$R$, le \omterm{def:b3:point-groups:representation}{caractère} de cette représentation de dimension $3N$ vaut
\[
\chi_{3N}(R) = (\text{nombre d'atomes laissés en place}) \times \chi_{xyz}(R),
\]
avec $\chi_{xyz} = 1 + 2\cos\theta$ pour une rotation d'angle $\theta$ et
$-1 + 2\cos\theta$
pour une rotation impropre (donc 3 pour $E$, $-1$ pour $C_2$, 0 pour
$C_3$, 1 pour $\sigma$,
$-3$ pour $i$).
\end{proposition}

\begin{proof}
Un atome envoyé sur une autre position emporte ses vecteurs hors de la
diagonale de la matrice~: il ne contribue pas à la trace. Un atome laissé
en place voit ses trois vecteurs transformés par la matrice $3\times3$ de
$R$, dont la trace se calcule dans un repère où $z$ est porté par
l'axe~: $\cos\theta + \cos\theta + 1$ pour une rotation, le dernier terme
devenant $-1$ pour une rotation impropre.
\end{proof}

\begin{proposition}[Dénombrer les vibrations]\label{prop:b3:group-theory-applied:mode-count}
$\Gamma_{\mathrm{vib}} = \Gamma_{3N} - \Gamma_{\mathrm{trans}} - \Gamma_{\mathrm{rot}}$, où
$\Gamma_{\mathrm{trans}}$ est la représentation de $(x, y, z)$ et $\Gamma_{\mathrm{rot}}$ celle
de $(R_x, R_y, R_z)$~; une molécule non linéaire a $3N - 6$ vibrations,
une molécule linéaire $3N - 5$.
\end{proposition}

\begin{proof}
Les $3N$ déplacements décrivent tous les mouvements~: trois translations
du centre de masse, trois rotations (deux pour une molécule linéaire,
puisque la rotation autour de son propre axe ne déplace aucun noyau), et
les vibrations. Ces sous-espaces ne se mélangent pas sous l'effet des
opérations, donc leurs représentations s'additionnent.
\end{proof}

\begin{example}[L'eau]\label{ex:b3:group-theory-applied:water}
% ledger: vib:H2O.sym, vib:H2O.bend, vib:H2O.asym
\ce{H2O} dans $C_{2v}$, molécule dans le plan $xz$. Atomes laissés en
place~: 3, 1, 3, 1~;
$\chi_{xyz}$~: $3, -1, 1, 1$~; donc $\chi_{3N} = 9, -1, 3, 1$. La
réduction donne $3A_1 +
A_2 + 3B_1 + 2B_2$. Une fois retirées les translations $A_1 + B_1 + B_2$
et les rotations $A_2 + B_1 + B_2$,
il reste $\Gamma_{\mathrm{vib}} = 2A_1 + B_1$. Les deux modes $a_1$ sont
l'élongation symétrique
(\qty{3657}{cm^{-1}}) et la déformation angulaire (\qty{1595}{cm^{-1}})~;
le mode $b_1$ est l'élongation antisymétrique (\qty{3756}{cm^{-1}}).
\end{example}

\begin{omfigure}
\begin{tikzpicture}[font=\small, >=Stealth]
  \foreach \x/\name/\lab in {0/sym/{$\nu_1$, $a_1$, élongation symétrique}, 4.3/bend/{$\nu_2$, $a_1$, déformation}, 8.6/asym/{$\nu_3$, $b_1$, élongation antisymétrique}}{
    \begin{scope}[xshift=\x cm]
      \node[aO] (o\name) at (0,0.35) {};
      \node[aH] (a\name) at (-0.8,-0.25) {};
      \node[aH] (b\name) at (0.8,-0.25) {};
      \begin{scope}[on background layer]\draw[bond] (o\name.center) -- (a\name.center) (o\name.center) -- (b\name.center);\end{scope}
      \node[anchor=north, align=center, font=\footnotesize] at (0,-0.9) {\lab};
    \end{scope}}
  % symmetric stretch: H outward along the bonds, O up a little
  \draw[->, omThm, thick] (-0.95,-0.36) -- (-1.42,-0.71);
  \draw[->, omThm, thick] (0.95,-0.36) -- (1.42,-0.71);
  \draw[->, omThm, thick] (0,0.74) -- (0,0.98);
  % bend: each H moves perpendicular to its bond, closing the angle; O up
  \draw[->, omThm, thick] (4.3-0.92,-0.36) -- (4.3-0.62,-0.76);
  \draw[->, omThm, thick] (4.3+0.92,-0.36) -- (4.3+0.62,-0.76);
  \draw[->, omThm, thick] (4.3,0.74) -- (4.3,0.98);
  % antisymmetric stretch: one H out, the other in, O sideways
  \draw[->, omThm, thick] (8.6-0.95,-0.36) -- (8.6-1.42,-0.71);
  \draw[->, omThm, thick] (8.6+1.02,0.0) -- (8.6+0.70,0.24);
  \draw[->, omThm, thick] (8.6-0.05,0.80) -- (8.6+0.25,0.80);
\end{tikzpicture}

{\small Les trois modes normaux de l'eau (flèches~: déplacements, pas à
l'échelle). Les deux modes $a_1$ conservent toute la symétrie de la
molécule~; le mode $b_1$ change de signe par $C_2$.}
\end{omfigure}

\begin{example}[Quatre autres molécules]\label{ex:b3:group-theory-applied:table}
La même procédure (calculée et vérifiée dans les données de figure de ce
chapitre) donne~:
\begin{center}\small
\begin{tabular}{llll}
\hline
molécule & groupe & $\Gamma_{\mathrm{vib}}$ & modes \\ \hline
\ce{NH3} & $C_{3v}$ & $2A_1 + 2E$ & 6 \\
\ce{CH4} & $T_d$ & $A_1 + E + 2T_2$ & 9 \\
\ce{CO2} & $D_{\infty h}$ (par $D_{2h}$) & $\Sigma_g^+ + \Sigma_u^+ + \Pi_u$ ($A_g + B_{1u} + B_{2u} + B_{3u}$) & 4 \\
\ce{BF3} & $D_{3h}$ & $A_1' + 2E' + A_2''$ & 6 \\
\ce{XeF4} & $D_{4h}$ & $A_{1g} + B_{1g} + B_{2g} + A_{2u} + B_{2u} + 2E_u$ & 9 \\
\ce{SF6} & $O_h$ & $A_{1g} + E_g + T_{2g} + 2T_{1u} + T_{2u}$ & 15 \\ \hline
\end{tabular}
\end{center}
Pour une molécule linéaire, on remplace le groupe infini par son
\omterm{def:b3:point-groups:class}{sous-groupe} $D_{2h}$,
dans lequel la déformation dégénérée $\Pi_u$ se scinde en $B_{2u} + B_{3u}$.
\end{example}

\section{Règles de sélection}

\begin{definition}[Produit direct]\label{def:b3:group-theory-applied:direct-product}
Le \emph{produit direct}\index{produit direct} $\Gamma_i\otimes\Gamma_j$
de deux représentations est la représentation portée par les produits de
leurs fonctions de base~; ses \omterm{def:b3:point-groups:representation}{caractères} sont les produits
$\chi_i(R)\chi_j(R)$.
\end{definition}

\begin{theorem}[Intégrales nulles]\label{thm:b3:group-theory-applied:vanishing-integral}
Une intégrale $\int f_1f_2f_3\,\dd\tau$ sur tout l'espace ne peut être
non nulle que si le
\omterm{def:b3:group-theory-applied:direct-product}{produit direct} $\Gamma_1\otimes\Gamma_2\otimes\Gamma_3$ contient la
\omterm{def:b3:group-theory-applied:reducible}{représentation totalement symétrique}.
\end{theorem}

\begin{proof}
Une \omterm{def:b3:point-groups:operation}{opération de symétrie} ne fait que renuméroter les points de l'espace,
donc la valeur de l'intégrale est inchangée quand l'intégrande est
transformée par une opération $R$ quelconque. Faisons la moyenne de
l'intégrande transformée sur le groupe~: la moyenne des composantes de
l'intégrande appartenant à chaque \omterm{def:b3:point-groups:irrep}{représentation irréductible} est
$\frac1h\sum_R\hat R(\cdot)$, qui est la projection sur la \omterm{def:b3:group-theory-applied:reducible}{représentation
totalement symétrique}
(\cref{def:b3:group-theory-applied:projection} avec tous les
$\chi = 1$). Toute autre composante a une moyenne nulle~; si le produit
ne contient aucune partie totalement symétrique, l'intégrale est nulle.
\end{proof}

\begin{definition}[Actif en IR, actif en Raman]\label{def:b3:group-theory-applied:activity}
Une vibration fondamentale est \emph{active en IR}\index{actif en IR} si
elle peut absorber le rayonnement infrarouge, \emph{active en Raman}\index{actif en Raman} si elle
apparaît dans le spectre Raman (\cref{ch:b3:rovibrational-spectroscopy}).
\end{definition}

\begin{corollary}[Règles de sélection infrarouge et Raman]\label{cor:b3:group-theory-applied:ir-raman}
Une transition fondamentale (du niveau fondamental, totalement
symétrique, à un quantum d'un mode de symétrie $\Gamma$) n'est \omterm{def:b3:group-theory-applied:activity}{active en
IR} que si $\Gamma$ est la
représentation de $x$, $y$ ou $z$, et n'est \omterm{def:b3:group-theory-applied:activity}{active en Raman} que si
$\Gamma$ est la représentation d'une fonction quadratique ($x^2$, $xy$,
\dots).
\end{corollary}

\begin{proof}
L'intensité de l'absorption fait intervenir $\int\psi_0\,\mu_k\,\psi_1\dd\tau$
avec les composantes du dipôle
$\mu_k$, qui se transforment comme $x$, $y$, $z$~; $\psi_0$ est
totalement symétrique
et $\psi_1$ se transforme comme $\Gamma$. D'après le théorème, l'intégrale
ne peut être non nulle
que si $\Gamma\otimes\Gamma_{\mu_k}$ contient $A_1$, ce qui n'arrive que si $\Gamma =
\Gamma_{\mu_k}$ (le produit de deux \omterm{def:b3:point-groups:irrep}{représentations irréductibles} ne
contient la \omterm{def:b3:group-theory-applied:reducible}{représentation totalement symétrique} que si elles sont
égales, pour des \omterm{def:b3:point-groups:representation}{caractères} réels). La diffusion Raman fait intervenir
les composantes $\alpha_{kl}$ de la polarisabilité, qui se transforment
comme les produits $kl$.
\end{proof}

\begin{proposition}[Règle d'exclusion mutuelle]\label{prop:b3:group-theory-applied:mutual-exclusion}
Dans une molécule possédant un \omterm{def:b3:point-groups:inversion}{centre d'inversion}, aucune transition
fondamentale n'est à la fois \omterm{def:b3:group-theory-applied:activity}{active en IR} et en Raman~: c'est la
\emph{règle d'exclusion mutuelle}\index{regle d'exclusion mutuelle@règle d'exclusion mutuelle}.
\end{proposition}

\begin{proof}
Avec un \omterm{def:b3:point-groups:inversion}{centre d'inversion}, toute \omterm{def:b3:point-groups:irrep}{représentation irréductible} est $g$ ou
$u$. Les coordonnées $x$, $y$, $z$ changent de signe par $i$ (elles sont
$u$)~; les fonctions quadratiques non ($g$). Un mode n'est \omterm{def:b3:group-theory-applied:activity}{actif en IR}
que s'il est $u$, \omterm{def:b3:group-theory-applied:activity}{actif en Raman} que s'il est
$g$.
\end{proof}

\begin{example}[Le dioxyde de carbone et l'effet de serre]\label{ex:b3:group-theory-applied:co2}
% ledger: vib:CO2.sym, vib:CO2.bend, vib:CO2.asym, ir:CO2.asym.int, ir:CO2.bend.int
\ce{CO2} possède un \omterm{def:b3:point-groups:inversion}{centre d'inversion}. Son élongation symétrique
($\Sigma_g^+$,
\qty{1333}{cm^{-1}}) n'est active qu'en Raman~; l'élongation
antisymétrique ($\Sigma_u^+$,
\qty{2349}{cm^{-1}}) et la déformation ($\Pi_u$, \qty{667}{cm^{-1}}) ne
sont actives qu'en IR.
La déformation absorbe vers \qty{15}{\micro m}, près du maximum de
l'émission thermique de la Terre~: c'est pourquoi le dioxyde de carbone
est un gaz à effet de serre. \ce{N2} et
\ce{O2} n'ont qu'une vibration, $\Sigma_g^+$~: inactive en IR.
\end{example}

\begin{omfigure}
\begin{tikzpicture}
\begin{axis}[omir, width=0.92\linewidth, height=4.4cm, xmin=400, xmax=2600,
  ymin=0, ymax=1.3, ytick=\empty, xlabel={nombre d'onde / cm$^{-1}$},
  ylabel={signal}, legend cell align=left,
  legend style={font=\small, at={(0.98,0.95)}, anchor=north east, draw=none}]
\addplot[omDef, thick] table[x=nu, y=ir] {figdata/out/bachelor-3-co2-spectra.dat};
\addlegendentry{absorption infrarouge}
\addplot[omThm, thick, dashed] table[x=nu, y=raman] {figdata/out/bachelor-3-co2-spectra.dat};
\addlegendentry{diffusion Raman}
\node[font=\small, anchor=south] at (axis cs:2349,1.02) {$\Sigma_u^+$};
\node[font=\small, anchor=south] at (axis cs:667,0.1) {$\Pi_u$};
\node[font=\small, anchor=south, omThm] at (axis cs:1333,1.02) {$\Sigma_g^+$};
\end{axis}
\end{tikzpicture}

{\small Spectres synthétiques de \ce{CO2} gazeux (positions des bandes et
intensités infrarouges relatives tirées des valeurs mesurées~; formes
schématiques, sans structure de rotation). Exclusion mutuelle~: aucune
bande n'apparaît dans les deux spectres.}
\end{omfigure}

\begin{method}[Prévoir un spectre de vibration]\label{met:b3:group-theory-applied:vibrations}
\begin{enumerate}
  \item Trouver le \omterm{def:b3:point-groups:point-group}{groupe ponctuel}~; calculer $\chi_{3N}$ à partir des
        atomes laissés en place.
  \item Réduire, retrancher les translations et les rotations~:
        $\Gamma_{\mathrm{vib}}$.
  \item À l'aide des colonnes de droite de la table, marquer chaque mode
        \omterm{def:b3:group-theory-applied:activity}{actif en IR}
        ($x$, $y$, $z$) et/ou \omterm{def:b3:group-theory-applied:activity}{actif en Raman} (quadratique)~; un mode qui
        n'est ni l'un ni l'autre est silencieux.
  \item Dénombrer les bandes~: une par mode actif, un mode dégénéré ne
        donnant qu'une bande.
\end{enumerate}
\end{method}

\begin{method}[Dénombrer les élongations CO]\label{met:b3:group-theory-applied:co-count}
Dans un métal carbonyle, prendre pour base les $n$ vecteurs des liaisons
C--O (un atome laissé en
place compte pour 1, rien d'autre ne compte)~: réduire pour obtenir
$\Gamma_{\mathrm{CO}}$, puis appliquer
les règles IR et Raman. Le nombre de bandes intenses vers
\qtyrange{1850}{2150}{cm^{-1}} identifie l'isomère.
\end{method}

\begin{example}[Cis et trans]\label{ex:b3:group-theory-applied:cis-trans}
cis-\ce{ML4(CO)2}, $C_{2v}$~: \omterm{def:b3:point-groups:representation}{caractères} $2, 0, 2, 0$, $\Gamma_{\mathrm{CO}} = A_1 + B_1$,
deux bandes IR. trans-\ce{ML4(CO)2}, $D_{4h}$~: $\Gamma_{\mathrm{CO}} = A_{1g} + A_{2u}$,
une bande IR ($A_{2u}$) et une bande Raman ($A_{1g}$). Une seule bande CO
en infrarouge signe l'isomère trans.
\end{example}

\section{Application aux transitions électroniques}

Le même théorème s'applique aux transitions électroniques~: une transition
électronique entre les états $\Psi_1$ et $\Psi_2$ n'est permise que si
$\Gamma_1\otimes\Gamma_2$
contient la représentation de $x$, $y$ ou $z$, et sa lumière est
polarisée selon cet axe.

\begin{example}[Transitions de l'eau et du méthanal]\label{ex:b3:group-theory-applied:electronic}
Dans l'eau ($C_{2v}$, plan $xz$), promouvoir un électron de $1b_2$ (le
doublet non liant) vers $4a_1$ donne un état excité de symétrie
$B_2\otimes A_1 = B_2$, qui est celle de $y$~:
transition permise, polarisée perpendiculairement au plan de la molécule.
Dans le méthanal
($C_{2v}$), la transition $n\to\pi^*$ part d'un doublet non liant de
l'oxygène situé dans le plan ($b_1$) vers l'orbitale $\pi^*$, construite
à partir d'orbitales $p$ perpendiculaires au plan
($b_2$~; molécule dans le plan $xz$, C=O selon $z$)~: $B_1\otimes B_2 = A_2$,
qui n'est la symétrie d'aucune des coordonnées $x$, $y$, $z$~: transition
interdite par la symétrie, et c'est pourquoi sa bande
(\cref{ch:b3:electronic-spectroscopy}) est si faible.
\end{example}

\begin{history}[Wigner, Bethe et la symétrie des molécules]
La théorie des groupes entra en mécanique quantique avec les travaux
d'Eugene Wigner sur les spectres atomiques (1927--31) et l'article de
Hans Bethe sur l'éclatement des termes atomiques dans les cristaux
(1929). E. Bright Wilson l'appliqua aux vibrations moléculaires en
1934~; Robert Mulliken donna les symboles encore utilisés pour les
orbitales et les états.
\end{history}

\begin{inthelab}[Infrarouge et Raman côte à côte]
Un spectromètre infrarouge mesure la lumière qu'absorbe un échantillon~;
un spectromètre Raman l'éclaire avec un laser et analyse la faible
lumière diffusée à des longueurs d'onde décalées. Enregistrer les deux
spectres d'un même composé est le test classique d'un centre de
symétrie~: si aucune bande ne coïncide, la molécule est probablement
centrosymétrique. Les lasers Raman sont de classe 3B ou 4~: le faisceau
est capoté et les opérateurs portent des lunettes adaptées à sa longueur
d'onde.
\end{inthelab}

\section{Exercices}

\begin{exercise}[$\star$]\label{exo:b3:group-theory-applied:1}
Réduire la représentation de $C_{2v}$ de \omterm{def:b3:point-groups:representation}{caractères} $(5, -1, 3, 1)$ sur $(E, C_2,
\sigma_v(xz), \sigma_v'(yz))$. Vérifier la dimension.
\end{exercise}

\begin{exercise}[$\star$]\label{exo:b3:group-theory-applied:2}
\ce{SO2} est coudée, comme l'eau. Donner $\Gamma_{\mathrm{vib}}$ et dire
quels modes sont \omterm{def:b3:group-theory-applied:activity}{actifs en IR} et en Raman.
\end{exercise}

\begin{exercise}[$\star$]\label{exo:b3:group-theory-applied:3}
Dans $C_{2v}$, calculer les \omterm{def:b3:group-theory-applied:direct-product}{produits directs} $B_1\otimes B_2$ et
$A_2\otimes B_1$, puis
$E\otimes E$ dans $C_{3v}$ (réduire ce dernier).
\end{exercise}

\begin{exercise}[$\star$]\label{exo:b3:group-theory-applied:4}
% ledger: vib:CH4.nu1, vib:CH4.nu2, vib:CH4.nu3, vib:CH4.nu4
\ce{CH4} a des transitions fondamentales à 2917 ($a_1$), 1534 ($e$), 3019
et 1306 (toutes deux $t_2$)
\unit{cm^{-1}}. Lesquelles apparaissent dans le spectre infrarouge,
lesquelles dans le spectre Raman~?
\end{exercise}

\begin{exercise}[$\star\star$]\label{exo:b3:group-theory-applied:5}
Pour \ce{BF3}, établir $\Gamma_{\mathrm{vib}} = A_1' + 2E' + A_2''$ à
partir des atomes laissés en
place, puis dénombrer les bandes IR et Raman. Quel mode n'apparaît que
dans un seul des deux spectres, et lequel dans les deux~?
\end{exercise}

\begin{exercise}[$\star\star$]\label{exo:b3:group-theory-applied:6}
Combien de bandes IR et Raman \ce{SF6} présente-t-il~? Quel mode est
silencieux~? Pourquoi aucune bande n'apparaît-elle dans les deux
spectres~?
\end{exercise}

\begin{exercise}[$\star\star$]\label{exo:b3:group-theory-applied:7}
Appliquer l'\omterm{def:b3:group-theory-applied:projection}{opérateur de projection} à $h_2$ pour la représentation $E$ de
$C_{3v}$ et
montrer qu'orthogonaliser le résultat par rapport à $2h_1 - h_2 - h_3$
donne $h_2 - h_3$.
\end{exercise}

\begin{exercise}[$\star\star$]\label{exo:b3:group-theory-applied:8}
Prévoir le nombre de bandes CO en IR du fac- et du mer-\ce{ML3(CO)3}, et
de \ce{M(CO)6}.
\end{exercise}

\begin{exercise}[$\star\star$]\label{exo:b3:group-theory-applied:9}
Montrer que les six orbitales $\sigma$ des ligands d'un complexe
octaédrique engendrent
$A_{1g} + E_g + T_{1u}$, et dire quelles orbitales du métal peuvent se
combiner avec chacune.
\end{exercise}

\begin{exercise}[$\star\star\star$]\label{exo:b3:group-theory-applied:10}
La représentation régulière a $\chi(E) = h$ et $\chi(R) = 0$ pour
$R \ne E$. Montrer
à l'aide de la \omterm{thm:b3:group-theory-applied:reduction}{formule de réduction} qu'elle contient chaque
\omterm{def:b3:point-groups:irrep}{représentation irréductible} $\Gamma_i$ exactement $d_i$ fois, et en
déduire $\sum_id_i^2 = h$.
\end{exercise}

\begin{exercise}[$\star\star\star$]\label{exo:b3:group-theory-applied:11}
L'acétylène \ce{HC#CH} est linéaire. En travaillant dans $D_{2h}$,
trouver $\Gamma_{\mathrm{vib}}$, le regrouper selon les symboles de
$D_{\infty h}$ et prévoir les spectres IR et Raman.
\end{exercise}

\begin{exercise}[$\star\star\star$]\label{exo:b3:group-theory-applied:12}
Dans l'eau (plan $xz$), lesquelles de ces promotions monoélectroniques
donnent des transitions permises, et avec quelle polarisation~:
$1b_2 \to 4a_1$, $3a_1 \to 4a_1$,
$1b_2 \to 2b_1$, $1b_1 \to 4a_1$~?
\end{exercise}

\section{Problème~: cis ou trans~? Un complexe carbonyle lu dans son spectre}

\begin{problem}[{Problème du week-end --- groupes ponctuels et
représentations des élongations CO de quatre isomères, leurs bandes
infrarouges et Raman, et l'identification d'un produit à partir de son
spectre mesuré}]
\label{pb:b3:group-theory-applied:1}
% ledger: ct:C2v, ct:C3v, ct:D4h
Un métal M octaédrique porte des ligands carbonyle et d'autres ligands L
(chaque L étant compté comme un point). On considère quatre isomères~:
le cis- et le
trans-\ce{ML4(CO)2}, le fac- et le mer-\ce{ML3(CO)3}. On suppose que les
ligands L n'abaissent pas la symétrie au-delà de ce qu'imposent leurs
positions. Une synthèse de
\ce{ML3(CO)3} donne un produit dont le spectre infrarouge présente trois
bandes intenses à 2048, 1965 et \qty{1938}{cm^{-1}} (données du
problème)~; son spectre Raman présente trois bandes à des positions
presque identiques.

\medskip\noindent\textbf{Partie I --- \omterm{def:b3:point-groups:point-group}{Groupes ponctuels}.}
\begin{enumerate}
  \item Dessiner les quatre isomères sur un octaèdre.
  \item Donner le \omterm{def:b3:point-groups:point-group}{groupe ponctuel} du cis-\ce{ML4(CO)2}.
  \item Du trans-\ce{ML4(CO)2}.
  \item Du fac-\ce{ML3(CO)3}.
  \item Du mer-\ce{ML3(CO)3}.
  \item Lesquels des quatre possèdent un \omterm{def:b3:point-groups:inversion}{centre d'inversion}~?
\end{enumerate}

\medskip\noindent\textbf{Partie II --- Les élongations CO.}
\begin{enumerate}[resume]
  \item Expliquer pourquoi les vecteurs des liaisons C--O portent une
        représentation dans laquelle une opération contribue pour 1 par
        CO laissé en place.
  \item Trouver $\Gamma_{\mathrm{CO}}$ de l'isomère cis.
  \item De l'isomère trans.
  \item De l'isomère fac.
  \item De l'isomère mer.
  \item Comparer chaque dimension au nombre de ligands CO.
\end{enumerate}

\medskip\noindent\textbf{Partie III --- Activité.}
\begin{enumerate}[resume]
  \item Dénombrer les élongations CO \omterm{def:b3:group-theory-applied:activity}{actives en IR} de chaque isomère.
  \item Dénombrer celles qui sont \omterm{def:b3:group-theory-applied:activity}{actives en Raman}.
  \item Quel isomère illustre la \omterm{prop:b3:group-theory-applied:mutual-exclusion}{règle d'exclusion mutuelle}~?
  \item Dans l'isomère trans, décrire le mode \omterm{def:b3:group-theory-applied:activity}{actif en IR}~: les deux CO
        vibrent-ils en phase ou en opposition de phase~?
  \item Dans l'isomère fac, pourquoi les trois CO ne donnent-ils que deux
        bandes~?
  \item Que révélerait un décompte de 1 bande (IR)~?
\end{enumerate}

\medskip\noindent\textbf{Partie IV --- Le produit.}
\begin{enumerate}[resume]
  \item Quel isomère le spectre infrarouge du produit désigne-t-il~?
  \item Le spectre Raman est-il cohérent~?
  \item La bande la plus haute est l'élongation en phase des groupes CO.
        À quelle \omterm{def:b3:point-groups:irrep}{représentation irréductible} appartient-elle dans cet
        isomère~?
  \item Le spectre pourrait-il provenir d'un mélange des deux isomères~?
        Quelle mesure supplémentaire trancherait~?
  \item Énoncer le résultat~: le nombre d'élongations CO \omterm{def:b3:group-theory-applied:activity}{actives en IR} de
        l'isomère mer, comparé à celui de l'isomère fac.
\end{enumerate}
\end{problem}
